% lägg in titlar

Let’s be honest: chances are you’re not going to read this entire thesis, and that’s totally fine. You’re probably here because this is the preface and it’s the perfect place to begin (and perhaps to end). We are all busy, right? Academics are busy working on papers (publish or perish!), and practitioners are swamped with back-to-back meetings and projects. But stick with me for just this short prologue.

Over the past five years, my PhD has kept me deeply interested in retail analytics. The questions I ask might sound familiar (as some might say, not particularly groundbreaking): Why do discounts work better in some stores than others? How does launching an online channel shift customer behavior?\footnote{Throughout this thesis, the terms consumer and customer are used with distinct but related meanings. The term consumer refers to general consumer behavior, particularly the psychological and decision-making processes often examined in consumer behavior research. The main focus of this thesis is on customers, taken as individuals actively interacting with retailers within the studied empirical settings. In other words, this thesis draws on consumer theory to explain how retailers might expect their customers to respond to strategic actions} What happens to customer behavior on digital platforms in the face of an external shock like COVID-19? These are questions that retailers, like consumers, ask all the time. What sets this work apart is my ability (and privilege) to spend half a decade or more exploring the questions with solid data, structured models, and consumer behavior theory to help connect the dots between observed patterns and strategic outcomes.
In this thesis, I look at how retailers adapt their pricing, channel, and platform strategies, and I examine how customers respond in ways that impact sales, profits, and outcomes on digital platforms. With the help of data and theory, I try to figure out what works, where, and why.
Here’s what you’ll discover.


\textbf{1. Multiformat: Pricing Strategy and Store Contexts}

Choosing a price is one thing, but it is another thing to know how customers respond differently to that price depending on where they shop. In the first two of the studies presented herein, I examine how pricing affects customer behavior across different physical store formats including hypermarkets, supermarkets, and convenience stores. The first study looks at private label tiers. Including basic, standard, and premium tiers, it shows that sales and profit elasticity vary significantly depending on the store format. The second study compares regular price reductions with promotional price discount. It shows that the same price cut can produce different effects depending on the context.

\textbf{2. Multichannel: The Long-Term Impact of Going Online}

In the third study, I investigate how online channel adoption changes customer behavior over the long run. The study finds that online adoption leads to long-term changes: larger basket sizes, increased spending in bulky or heavy categories like bottled water, and reduced purchases in impulse-driven or perishable items like snacks or fresh vegetables. These shifts reflect stable changes in how customers allocate effort and convenience across channels. Rather than cannibalizing in-store sales, online channels can enhance overall customer value.

\textbf{3. Platform: Customer Behavior During External Shocks}

In the final study, I examine how customer behavior on digital platforms shifted during the COVID-19 pandemic. The study uses follower growth on Spotify playlists as an empirical example to capture changes in this digital platform context. The results show that playlists curated by the platform remained resilient, but those associated with major record labels and superstar tracks lost momentum. The takeaway addresses more than music. In algorithmic environments where attention is scarce, resilience comes from diversified offerings, mid-tail strategies, and institutional curation, not from a sole product’s power or influence. For retailers operating on digital platforms, this highlights the importance (especially during periods of uncertainty) of visibility management and content strategy.

If you continue flipping through the pages of this thesis, you will come across theoretical discussions, mathematical notations, and statistical models. Such material might not be everyone’s cup of tea, but that is what academics thrive on: building on past research, applying theories to explain real-world phenomena, and using mathematical formulas as precision tools to quantify observations and draw conclusions. Don’t worry; the theoretical discussions and equations are here to shed light on how retail strategies, in pricing and platforms, connect with customer behavior. It is like building a house, where theory provides the blueprint, data supplies the raw materials, and models are the tools that assemble everything into something meaningful and durable.

This thesis is not a definitive guide to retail strategy. What I offer is a series of observations, tests, and tentative conclusions grounded in empirical data. The value is not just in the answers but in the act of asking, of looking closely, questioning assumptions, explaining patterns. I hope to show that doing research is less about undeniable truths and more about honest attempts. Attempts to understand, to explain, to make sense of a world that is constantly shifting, sometimes so dramatically that it breaks whatever model was in development. And in the end, maybe that’s the real strategy: stay curious, stay skeptical, and never trust a finding that required three robustness checks and a prayer.


%(\cite{Altmejd2018_predicting_replicability})


%\bigskip\fancybreak{{* * *}}\bigskip

%Some more really great text about the second paper, based on %\textcite{Camerer2016_evaluating_replicability}.
